\documentclass[11pt,a4paper]{article}
\usepackage[utf8]{inputenc}
\usepackage[T1]{fontenc}
\usepackage{amsfonts}
\usepackage[french]{babel}
\frenchbsetup{StandardLists=true} % à inclure si on utilise \usepackage[francais]{babel}
\usepackage{enumitem}
\usepackage{amssymb}
\usepackage{amsmath}
\DeclareMathOperator{\e}{e}
\usepackage[left=2cm,right=2cm,top=2cm,bottom=2cm]{geometry}
\usepackage{multirow}
\usepackage[dvipsnames]{xcolor}

\usepackage[T1]{fontenc}
\usepackage{fouriernc}
\usepackage{graphicx}

\usepackage{setspace}
\singlespacing

\usepackage{multicol}

\usepackage{fancyhdr}
\pagestyle{fancy}
\renewcommand\headrulewidth{1pt}
\fancyhead[L]{Lycée Denis Diderot }
\fancyhead[R]{Classe 1BTS}
\setcounter{page}{1}\fancyfoot[C]{page \thepage}
\fancyfoot[R]{\today}
\fancyfoot[L]{Alexis RUIZ}
\usepackage{multirow,tabularx,booktabs}
\newcommand{\cadreblanc}[1]{%
\medskip\framebox[\linewidth][c]{%
\rule{0mm}{#1mm}}\par
\medskip}

\setlength{\columnseprule}{.25mm}
\renewcommand{\arraystretch}{2}

\begin{document}
\center \LARGE{CONTRÔLE - Primitives 1  (45 minutes)}\
\\[0,3cm]

\normalsize

\hrule\
\\[0,2cm]

\flushleft \textbf{NOM :\\
Prénom:\\[0,2cm]
}
\hrule\
\\[0,2cm]






 
 \fcolorbox[gray]{0}{0.9}{\textbf{{Exercice 1. Calcul de primitives ~~(3 points) }}}\\[2pt]
 
 Donner une primitive des fonctions suivantes sur $I$. 
 
 \begin{multicols}{6}
\begin{center}
{
$f_1(x)=5$ \\
\cadreblanc {10}
$f_2(x)=x^3$\\
\cadreblanc {10}
$f_3(x)=-\dfrac{1}{x^2}$\\
\cadreblanc {10}
$f_4(x)=\dfrac{6}{x}$\\
\cadreblanc {10}
$f_5(x)=\e^x$\\
\cadreblanc {10}
$f_6(x)=2\cos(2x)$\\
\cadreblanc {10}
}
\end{center}
\end{multicols}

\fcolorbox[gray]{0}{0.9}{\textbf{{Exercice 2. Calcul de primitives ~~(4 points) }}}\\[2pt]
 Donner une primitive des fonctions suivantes sur $I$. 
\begin{multicols}{3}
  \begin{center}
  {
  $f_1 = 2x^2-6x+3$\\
  \cadreblanc {25}
  $f_2 = x+1+\dfrac{1}{x}$
  \cadreblanc{25}
  $f_3 = 2x(x^2+1)$\\
  \cadreblanc {25}
  }
  \end{center}
\end{multicols}

\fcolorbox[gray]{0}{0.9}{\textbf{{Exercice 3. Calcul d'une primitive ~~(2 points) }}}\\[2pt]

Soit $f$ la fonction définie sur $]0;+\infty[$ par $f(x)=2x-1$\\
Déterminer la primitive $F$ de $f$ sur $]0;+\infty[$ qui s'annule pour $x=2$.\\
\cadreblanc{40}

\fcolorbox[gray]{0}{0.9}{\textbf{{Exercice 4. La primitive est donnée ~~(1 points) }}}\\[2pt]
Vérifier que la fonction $F(x)=\dfrac{x+1}{x+2}$ est une primitive de la fonction $f(x)=\dfrac{1}{(x+2)^2}$.
\\
\cadreblanc {50}


\end{document}